% =====================================================================================
%  Sec. 1 -- THE OPERATIONAL QUESTION.  Replaces paper/introduction.tex in full.
%
%  WRITTEN LAST, against the numbers that came out.  Every quantity quoted here is
%  quoted somewhere else in the manuscript with its provenance; this section adds no
%  measurement of its own.  Measured length: 2.1 pages in the real preamble.
%
%  EXTERNAL LABELS USED (all owned by other sections; reconcile at integration):
%    sec:method     Sec. 2, the primitive; what is sampled and what is exact [exists, method.tex]
%    sec:leakage    Sec. 3, the two budgets / the leakage bound              [Sec. 3, not yet written]
%    thm:leakage    the two-sided bound of Sec. 3                           [Sec. 3, not yet written]
%    sec:retraction Sec. 3.1, retraction of Theorem 1(iii)                  [sec_3_1.tex]
%    sec:moments    Sec. 4, moment exactness                                [sec_4.tex]
%    thm:moments, prop:shots                                                [sec_4.tex]
%    sec:frontier   Sec. 5, the non-vacuity threshold, measured             [Sec. 5, not yet written]
%    sec:nogo       Sec. 6, the obstruction and its witness                 [sec_6.tex]
%    sec:nogo:chi   Sec. 6.5, the statistics the right way up               [sec_6.tex]
%    thm:witness, prop:leakinv, eq:collapse                                 [sec_6.tex]
%    sec:prices     Sec. 7, the three prices                                [sec_7.tex]
%    sec:physics    Sec. 8, the two gaps                                    [sec_8.tex]
%    tab:gaps                                                               [sec_8.tex]
%    sec:honest     Sec. 9, hardware, scope, priority                       [exists, discussion.tex]
%    fig:thm1iii-violation  new Fig. N3, Appendix B                         [p2_figs, caption written]
%    fig:gapscaling         new Fig. N2                                     [p2_figs, caption written]
%  This file defines only: sec:intro.
%
%  *** SEVEN \cite KEYS USED HERE DO NOT EXIST IN paper/bibliography.tex ***
%  They are real papers whose identifier fields are verified in this project's
%  state-of-the-art review; the TITLES are not verified here and must be pulled from
%  arXiv/Crossref before the \bibitem entries are written.  Do not guess a title.
%    \cite{patel2026}       Patel, Rangi, Tam, arXiv:2609.09147 (8 Sep 2026)
%    \cite{benthien2007}    Benthien, Jeckelmann, PRB 75, 205128 (2007),
%                           DOI 10.1103/PhysRevB.75.205128   [ALSO requested by sec_8.tex]
%    \cite{vilchez2025}     Vilchez-Estevez, Santos, Wang, Gambetta, PRB 112, 045143 (2025),
%                           DOI 10.1103/ydfw-k83n
%    \cite{granet2026akw}   Granet et al. (Quantinuum), arXiv:2605.01440 (2026)
%    \cite{lee2026sqw}      Lee, ..., Kandala, Banerjee, arXiv:2603.15608 (2026)
%    \cite{millar2026}      Millar et al., arXiv:2607.02673 (2026)
%    \cite{leonebittel2026} Leone, Bittel, arXiv:2607.29670 (31 Jul 2026)
%                           [ALSO requested by sec_6.tex.  Do NOT reuse \cite{bittel2025},
%                            which is Bittel-Mele-Eisert-Leone, arXiv:2409.17953.]
%  Existing keys reused here, all verified present in bibliography.tex:
%    quantinuum_dsf2026 (= Granet et al., arXiv:2607.07138), tarabunga2026, nocera2016.
%
%  GREP NOTES for the mechanical audit of this file:
%   (a) The arXiv identifier of this preprint appears NOWHERE in this file.  It is named
%       once in the manuscript, in Sec. 3.1, and once is the rule.
%   (b) The banned unity argument -- that one half of the paper could not be published on
%       its own -- is NOT made here in any form, and must not be reintroduced: it is a fact
%       about the arXiv calendar, not about the paper.  The join between the two halves is
%       made only where it is measured: the "+" of the bound, evaluated on S* = supp(phi).
%   (c) The words "decoupled" and "decoupling" appear nowhere in this file.  The measured
%       relation between F_1 and the support is a deterministic monotone one with an
%       unstable sign (Sec. 6.5), which is a different statement.
%   (d) No claim of quantum advantage is made; the explicit disclaimer is the seventh
%       paragraph and must survive copy-editing.
%   (e) The word "sampled" is used only where a budget is declared or where the object
%       is named as the infinite-shot limit of the declared protocol (Sec. 7).  The
%       fraction of the sector is the "published fraction" everywhere -- the name of
%       Sec. 5's table and of Sec. 7.5, which forbids calling it sampled without a shot
%       count.  "sampled fraction" was still used in seven places (the abstract among
%       them) until 2026-09-19; it is now used in none.  Generic uses say "sector
%       fraction", the phrase Table XI already carried.
% =====================================================================================

\section{The operational question}
\label{sec:intro}

A sample-based quantum calculation returns a list of bitstrings. Diagonalizing the Hamiltonian inside
the span of the configurations they name yields an energy, and---as this work and several others
show---a frequency-resolved response as well. What it does not yield is a statement of how wrong that
response is. This paper addresses the question a user is left with once the shots are spent:
\emph{given the subspace those shots populated, how wrong is the spectral function built on it?} The
answer has a positive part, statements that hold on any subspace, and a negative part, a measurement
showing that the resulting certificate is not informative at the sector fractions at which sampling
operates.

The positive part consists of three results. First, a two-sided a posteriori bound
(Theorem~\ref{thm:leakage}): for Rayleigh--Ritz reconstruction on any orthogonal projector, the $L_1$
error of the broadened spectral function lies between the Born weight of the probe that the subspace
failed to capture and that weight plus the Hamiltonian coupling that leaks across the boundary of the
subspace at the resolution $\eta$, a quantity with a closed form in the Ritz pairs the reconstruction
already produces. Second, the captured weight alone cannot play that role: an analytic counterexample
forces the constant of any bound indexed on it to its trivial value, and on the probe's own Born
support, where the captured weight is exactly one, the relative error is still $0.43$ at $L=6$ and
$0.35$ at $L=8$ at $\eta=0.18\,t$ on the periodic rings ($0.28$ on the open chain of
Sec.~\ref{sec:moments}; the two lattices are never pooled). Third, what the subspace structure does
guarantee is moment exactness: a subspace containing the order-$K$ Krylov space reproduces the spectral
moments through order $2K+1$ and no further (Theorem~\ref{thm:moments}), and a non-circular count gives
the shots that capture every high-probability configuration (Proposition~\ref{prop:shots}).

The negative part is measured. Evaluated on the subspaces of the resource scan, at five sizes up to a
sector of dimension $10\,306\,296$, the leakage branch of the bound is worse than the trivial bound at
every operating fraction of the resource scan, by factors of $1.66$ to $8.61$, and at $L=12$ that
verdict is converged in Krylov depth. The reconstructions at those fractions are therefore, at
present, uncertified; the one drawn at $85\%$ of its sector, Fig.~\ref{fig:akwsampled}, is certified,
loosely (Sec.~\ref{sec:frontier}). The reconstructions themselves are not in question: an independent
implementation recovers the acceptance fractions of the scan to under $2\%$. Two further measurements
locate what remains. The crossing of the trivial bound occurs at a true relative error between
$1.2\times10^{-3}$ and $7.3\times10^{-3}$ over sixteen $(L,\eta)$ cells whose true errors span five
decades, so the certificate is a calibrated a posteriori risk tracker, although a two-constant rule in
the missed weight, fitted out of sample, localizes the same crossing more tightly on eleven of fourteen
splits. And in the determinant basis in which the device measures, the Krylov space of the probe has
coordinate support on every determinant that the reflection about the seed site allows, measured from
$L=6$ to $L=14$, so at full captured weight only a subspace holding all of them has zero leakage; a
momentum-resolved probe confines that support to one momentum block, $1/L$ of the sector
(Sec.~\ref{sec:frontier:obstruction}).

The second half of the paper asks whether the certificate can be replaced by an a priori predictor
computed from the state before any circuit runs. The natural candidates are the efficiently
measurable monotones of fermionic non-Gaussianity, and on number-conserving states the one-body member
collapses to the single invariant $\mathcal F_1=4\,\mathrm{tr}[\gamma(1-\gamma)]$
(Eq.~\eqref{eq:collapse}). An invariant of the occupation spectrum cannot see a boundary: along a
geminal orbit every order of it is fixed while the determinant support runs from $O(K)$ to $2^{K}$ at
bond dimension $2$ (Theorem~\ref{thm:witness}), and on fibers where every invariant and the captured
weight are frozen the subspace-aware certificate still separates a non-vacuous value from the trivial
bound (Proposition~\ref{prop:leakinv}). Within the six $(L,N)$ strata of the Hubbard interaction sweep
the rank correlation of $\mathcal F_1$ with the determinant support is exactly $-1$ (exact permutation
$p=3.3\times10^{-13}$ against independent orderings), a deterministic monotone relation whose sign
reverses between families, as an orbital invariant compared with a basis-dependent count must
(Sec.~\ref{sec:nogo}). Sections~\ref{sec:prices} and~\ref{sec:physics} then price the method in
support, resolution and shots, and check the full-sector references against the opposite finite-size
scaling of the charge and spin gaps of the Hubbard chain.

Four bodies of work bound what is new here. \emph{(i)}~Patel, Rangi and Tam~\cite{patel2026}
construct Green's functions from sampled Krylov subspaces for dynamical mean-field
theory~\cite{dmft1996}, independently and concurrently, and cite this work as ``a closely related
sample-based approach to dynamical spectral functions, constructed directly from bitstring-sampled
subspaces''; we make no priority claim against it. \emph{(ii)}~Classical dynamical DMRG arrived
first: Benthien and Jeckelmann~\cite{benthien2007} computed all three momentum-resolved channels
reported here, for the half-filled extended Hubbard chain at $L=90$--$120$, nineteen years ago, and
correction-vector DMRG treats the Hubbard chain at $L=48$, the Heisenberg chain at $L=64$ and the
$t$--$J$ ladder at $48\times2$ sites~\cite{nocera2016}. \emph{(iii)}~Frequency-resolved responses of
correlated lattice models have been measured on devices far larger than ours: a retarded spin Green's
function of the Fermi--Hubbard chain at up to $30$ qubits~\cite{vilchez2025}, an
angle-resolved-photoemission-like $A(k,\omega)$ on a $27$-site chain at $54$
qubits~\cite{granet2026akw}, $S(q,\omega)$ benchmarked against a copper-sulfate
crystal~\cite{quantinuum_dsf2026} and on $50$ qubits against neutron scattering in
$\mathrm{KCuF_3}$~\cite{lee2026sqw}, and quench spectroscopy of $L=101$ spins~\cite{millar2026}. Our
device run spans its entire sector, so its agreement with exact diagonalization is exact by coverage
rather than a test of the method (Sec.~\ref{sec:honest}). \emph{(iv)}~The division between
Gaussian-invariant and basis-dependent quantities is due to Tarabunga \emph{et al.}~\cite{tarabunga2026},
six weeks before this work first appeared, and to Leone and Bittel~\cite{leonebittel2026}, seventeen
days before; the priority for it is theirs, and what this work adds is a constructive instance and its
extension from the determinant support to the leakage. \emph{We claim no quantum advantage in this
paper}: the one quantitative statement it makes about its own cost---that the basis-dependent bond
dimension tracks the determinant support---argues, wherever that bond dimension is small, for a tensor
network rather than for a sampler (Sec.~\ref{sec:honest}).

The paper is organized as follows. Section~\ref{sec:method} fixes the sampling primitive and states
which reconstructions carry a shot budget and which are full-sector references. Section~\ref{sec:leakage}
proves the bound and shows that no bound on the captured weight alone can replace it;
Sec.~\ref{sec:moments} states what the subspace structure certifies; Sec.~\ref{sec:frontier} evaluates
the certificate; Sec.~\ref{sec:nogo} addresses a priori predictors; Secs.~\ref{sec:prices}
and~\ref{sec:physics} price the method and validate the references; and Sec.~\ref{sec:honest} reports
the device runs, the molecular suite and the region without advantage. The systems are the half-filled
one-dimensional Hubbard model at $U/t=8$, rings and open chains at $L=4$ to $14$, and nineteen molecular
active spaces. Three different counts share the symbol $\mathcal S$ and are kept apart: the static
ground-state support of Sec.~\ref{sec:nogo}, the dynamical support of Sec.~\ref{sec:prices}, and the set
of configurations the sampler returned, the $\mathcal S$ of Theorem~\ref{thm:leakage}. Proofs, full
tables and extended discussions are in the Supplemental Material, cited as Sec.~S$n$.
